% Generated draft: requires mathematical and source review.
\documentclass[11pt]{article}
\usepackage[margin=25mm]{geometry}
\usepackage{fontspec}
\setmainfont{Latin Modern Roman}
\usepackage{amsmath,amssymb,mathtools,unicode-math}
\setmathfont{Latin Modern Math}
\usepackage{enumitem,needspace,xurl,hyperref}
\setlength{\parindent}{0pt}
\setlength{\parskip}{5pt}
\setlength{\emergencystretch}{3em}
\newcommand{\N}{\mathbb{N}}
\newcommand{\Z}{\mathbb{Z}}
\newcommand{\Q}{\mathbb{Q}}
\newcommand{\R}{\mathbb{R}}
\newcommand{\C}{\mathbb{C}}
\newcommand{\F}{\mathbb{F}}
\newcommand{\A}{\mathbb{A}}
\newcommand{\PP}{\mathbb{P}}
\newcommand{\E}{\mathbb{E}}
\newcommand{\eps}{\varepsilon}
\newcommand{\dd}{\,\mathrm{d}}
\begin{document}
\section*{UnsolvedMath ID: 30002590}
\textbf{Source problem code:} catalogue entry
\subsection{Definitions and mathematical statement}
The supplied question is
\[
\text{Does }\operatorname{hdet}(\mu_A)=1\text{ hold for every Artin--Schelter Gorenstein algebra }A?
\]
Here \(\mu_A\) denotes the Nakayama automorphism and \(\operatorname{hdet}\) denotes the homological determinant. The intended class of algebras and the hypotheses ensuring that these objects exist are not specified in the supplied record, so the unrestricted quantifier cannot yet be given a fully precise domain.

For terminology, in a standard connected graded setting one starts with a field \(k\) and an algebra \(A=\bigoplus_{n\geq0}A_n\) with \(A_0=k\), finite-dimensional graded pieces, and augmentation module \(k=A/A_{\geq1}\). The Artin--Schelter Gorenstein condition requires finite left and right injective dimensions equal to an integer \(d\), together with graded Ext conditions
\[
\underline{\operatorname{Ext}}^i_A(k,A)\cong
\begin{cases}
0,&i\ne d,\\
k(\ell),&i=d,
\end{cases}
\qquad
\underline{\operatorname{Ext}}^i_{A^{\mathrm{op}}}(k,A)\cong
\begin{cases}
0,&i\ne d,\\
k(\ell),&i=d,
\end{cases}
\]
for a grading shift \(\ell\); the first is an isomorphism of graded right modules and the second of graded left modules. This describes a standard convention, not additional hypotheses supplied for the original question. The source does not specify its base field, grading convention, or whether it assumes connectedness, local finiteness, or noetherianity.

A Nakayama automorphism describes the automorphism twist in an appropriate dualizing bimodule or complex, when such a description exists. In the graded setting, the homological determinant is a scalar character of suitable graded automorphisms arising from their action on the one-dimensional Gorenstein duality data. Its precise construction and scalar convention must be fixed. In particular, the intended hypotheses must ensure that \(\mu_A\) exists and belongs to the domain of \(\operatorname{hdet}\). The displayed question preserves the supplied target without silently restricting it to the illustrative connected graded setting.

The supplied dated assessment reports the equality for noetherian connected Artin--Schelter Gorenstein algebras satisfying an additional condition written \(D\epsilon\ne0\), including all noetherian Artin--Schelter regular cases. It reports that automatic nonvanishing and the unrestricted conjecture remain open. Neither \(D\) nor \(\epsilon\) is defined in the record, so this reported partial result cannot be restated here as a fully specified theorem. These are supplied literature qualifications, not independently verified mathematical or status claims.
\subsection{Short English statement}
The question asks whether the Nakayama automorphism of every Artin--Schelter Gorenstein algebra has homological determinant \(1\). The intended algebra class and the hypotheses defining these invariants need clarification. Supplied dated curation labels the problem partially solved and reports a result under additional assumptions, including an undefined condition \(D\epsilon\ne0\).
\subsection{Sources}
\begin{itemize}
\item Local UnsolvedMath registry, permanent ID \(30002590\), source code \texttt{catalogue entry}, title ``Homological Determinant of the Nakayama Automorphism,'' in the Oberwolfach Reports Open Problems collection. No individual proposer is supplied; the proposed year is recorded as 2014.
\item Supplied workshop provenance: \emph{Interactions between Algebraic Geometry and Noncommutative Algebra} (2015), Oberwolfach Reports, pp.\ 1365--1402, DOI \texttt{10.4171/owr/2014/25}. \url{https://doi.org/10.4171/owr/2014/25}. The supplied background describes these reports as records of workshops and their open-problem discussions.
\item Supplied literature reference: Reyes--Rogalski--Zhang, ``skew Calabi--Yau categories.'' \url{https://doi.org/10.1090/tran/6640}.
\item Supplied literature reference: Reyes--Rogalski--Zhang, arXiv version. \url{https://arxiv.org/abs/1408.0536}.
\item Literature assessment date: \texttt{2026-08-21}. The partially solved label, the reported result under \(D\epsilon\ne0\), its stated inclusion of all noetherian Artin--Schelter regular cases, and the remaining-open claims are supplied dated curation derived from the OWR record, not independently verified here. The source requests independent verification of time-sensitive status and cited claims.
\item The registry labels the statement \texttt{corrected\_verified} and says that notation, typography, and wording were normalized without changing mathematical content. This is supplied registry metadata, not independent verification of the formulation or a mathematical review.
\end{itemize}
\end{document}
